% This syllabus template was created by:
% Brian R. Hall
% Assistant Professor, Champlain College
% www.brianrhall.net

% Document settings
\documentclass[11pt]{article}
\usepackage[margin=1in]{geometry}
\usepackage[pdftex]{graphicx}
\usepackage{multirow}
\usepackage{setspace}
\usepackage{listings}
\pagestyle{plain}
\setlength\parindent{0pt}

\begin{document}
	
	% Course information
	\begin{tabular}{ l l }
		
		& \LARGE Physique-chimie \\\\
		& \LARGE Printemps 2026 \\\\
	\end{tabular}
	\vspace{10mm}
	
	% Course Outline
	
	Le contenu peut être modifié en fonction des progrès des étudiants.
	
	Semaine 2  
	
	\subsection*{Rappel de la notion de pression. Travail exercé par le gaz ou sur le gaz}
	
	La pression est définie comme la force divisée par la surface :
	\begin{lstlisting}
		P = F/S
	\end{lstlisting}
	
	Le travail est la force multipliée par le déplacement de son point d'application :
	\[
	W = P \times S \times dx
	\]
	
	Le travail sur le gaz peut être positif ou négatif. Dans le second cas, on dit que le gaz exerce le travail.
	\begin{itemize}
		\item le gaz est comprimé : il exerce un travail négatif ;
		\item le gaz se détend : il exerce un travail positif.
	\end{itemize}
	
	On modifie l'énergie d'un système lorsqu'on exerce un travail.
	
	\subsection*{L'autre moyen de changer l'énergie d'un système : le transfert de chaleur}
	
	On modifie l'énergie d'un système lorsqu'on lui fournit ou qu'on lui retire de la chaleur.
	
	Dans les processus réversibles, la quantité de chaleur échangée avec le système est définie par le changement d'un paramètre : l'entropie :
	\[
	Q = T \times (S_{apres}- S_{avant})
	\]
	
	où $Q$ est la quantité de chaleur échangée avec le système, $T$ la température absolue et $S$ l'entropie du système.
	
	C'est à la fois une définition de l'entropie et de la température absolue.
	
	Si le processus est irréversible, l'entropie peut augmenter même en l'absence d'échange de chaleur.
	
	\subsection*{Le bilan d’énergie}
	
	\[
	\Delta U = W + Q
	\]
	
	Pour un processus réversible :
	\[
	\Delta U = W + T \times \Delta S 
	\]
	
	\subsubsection*{Les cas particuliers}
	
	\begin{itemize}
		\item À volume constant : il n'y a pas de travail mécanique. Toute la chaleur sert à modifier l'énergie interne en changeant la température.
		\[
		\Delta U = Q
		\]
		On introduit la capacité thermique à volume constant $C_V$ :
		
		c'est la chaleur fournie pour une variation de température à volume constant :
		\[
		Q = C_V \,\Delta T
		\]
		
		\item À pression constante : il y a un travail mécanique sur le gaz.
		$
		\Delta U = Q - P\,\Delta V
	$
		On introduit la capacité thermique à pression constante $C_P$ :
	$
		Q = C_P \,\Delta T
$
		
		On peut aussi remarquer que
\[	Q = \Delta U + P\,\Delta V	\]
	

		
		ou encore
	
		\[ Q = \Delta H \]  avec  \[ H = U + PV\]
	
		
		$H$ est appelée l'enthalpie. À pression constante, la quantité de chaleur nécessaire pour changer l'état d'un corps (fusion, évaporation) correspond à la variation de son enthalpie ; on parle alors d'enthalpie de fusion ou d'évaporation.
	\end{itemize}
	
	La capacité thermique peut être massique (par kilogramme) ou molaire (par mole).
	
	\subsection*{Modèle du gaz parfait et équation d’état}
	
	\[
	PV = nRT
	\]
	
	$n$ : nombre de moles, \quad $R = 8{,}3 \ \mathrm{J \cdot mol^{-1} \cdot K^{-1}}$
	
	Transformations : adiabatique, isotherme, isochore.
	
	Représentation graphique sur les diagrammes :
	\begin{itemize}
		\item $(P, V)$
		\includegraphics[width=0.4\linewidth]{diagramme_pv_isobare_isotherme_isochore.jpg}
		\item $(T, V)$
	\end{itemize}
	
	La capacité thermique à volume constant d’un gaz parfait dépend du nombre d’atomes dans ses molécules : plus il est grand, plus la capacité est élevée. Pour un gaz monoatomique, sa capacité thermique molaire est :
	\[
	C_V = \frac{3}{2} R
	\]
	\newpage
	\subsection*{Le gaz réel et ses diagrammes} 
 Diagramme de phase de Clapeyron
 
 
		\includegraphics[width=0.7\linewidth]{courbe_rosee_isothermes.jpg}
 
Point critique pour eau:
T = 374°C , P=218 atmosphères (221 bar)

\subsection*{A venir} 
Le diagramme de Mollier représente plusieurs caractéristiques du système liquide-vapeur. Il sera étudié à la leçon suivante.
	
\end{document}



